\section{A core calculus for \VG{} nodes}\label{sec:calculus}

This section fixes the objects the later theorems are about. The calculus covers the MVP
fragment of the specification: scalar, collection and float operations, the bounded
combinators \op{SIMD\_MAP}, \op{FOLD\_BOUNDED}, \op{GENERATE} and \op{ASSERT\_GUARD}, calls by
hash, and three effect operations. It leaves out what the MVP rejects: recursion, sum types,
generics and explicit memory regions. Individual operations are given by a table
(\cref{app:ops}); the calculus needs only their typing, domain and meaning.

\subsection{Types and values}

\begin{definition}[Types]\label{def:types}
\[
\begin{array}{rcl}
\tau & ::= & \mathsf{Bool} \mid \mathsf{u}N \mid \mathsf{i}N \mid \mathsf{f16} \mid
  \mathsf{bf16} \mid \mathsf{f32} \mid \mathsf{f64}
  \mid \mathsf{Array}\langle\tau, n\rangle \mid \mathsf{Tensor}\langle\tau,[d_1,\dots,d_k]\rangle
  \mid \mathsf{Slice}\langle\tau\rangle \\
  & \mid & \mathsf{Tuple}\langle\tau_1,\dots,\tau_k\rangle \mid
  \mathsf{Result}\langle\tau,\tau'\rangle \mid \Eff \mid \mathsf{FileIn} \mid \mathsf{FileOut}
  \mid \mathsf{Clock}
\end{array}
\]
with $N\in\{8,16,32,64\}$ and static $n, d_i$. A type is \emph{fixed-size} if it is
$\mathsf{Bool}$, a scalar, or an array, tensor, tuple or result built from fixed-size types.
\end{definition}

$\sem{\mathsf{u}N} = \{0,\dots,2^N-1\}$ and $\sem{\mathsf{i}N}=\{-2^{N-1},\dots,2^{N-1}-1\}$;
arrays, tensors and slices denote sequences (a slice of length below $2^{64}$), tuples and
results the usual products and tagged sums, $\sem{\Eff}=\{\star\}$, and handle types denote
opaque capabilities that only the runtime creates.

\begin{definition}[Float values with one NaN]\label{def:onenan}
Let $\varphi$ be one of the binary formats binary16, bfloat16, binary32 or binary64, of width
$w_\varphi$ bits. Let $\mathrm{Fin}_\varphi$ be its finite values, with $+0$ and $-0$ distinct,
and $V_\varphi = \mathrm{Fin}_\varphi \cup \{+\infty,-\infty,\NaN\}$. Let
$B_\varphi = \{0,1\}^{w_\varphi}$ be the bit patterns. The decoding
$\dec_\varphi : B_\varphi\to V_\varphi$ is the IEEE~754 decoding~\cite{ieee754-2019}, except
that every pattern with an all-ones exponent field and a non-zero significand decodes to the
single value $\NaN$. The encoding $\enc_\varphi : V_\varphi\to B_\varphi$ maps every non-NaN
value to its IEEE~754 encoding and $\NaN$ to the canonical quiet NaN $c_\varphi$ with a zero
sign bit: \texttt{0x7E00}, \texttt{0x7FC0}, \texttt{0x7FC00000} and
\texttt{0x7FF8000000000000} for the four formats. We write $\sem{\mathsf{f}N}=V_\varphi$.
\end{definition}

This is the float domain of the SMT-LIB theory of floating point~\cite{brain2015fp}, which has
a single NaN, and the NaN that the deterministic profile of WebAssembly
produces~\cite{wasm-spec}.

\begin{lemma}[Encoding and decoding]\label{lem:encdec}
$\dec_\varphi\circ\enc_\varphi=\mathrm{id}_{V_\varphi}$, and $\enc_\varphi\circ\dec_\varphi$ is
idempotent and fixes exactly the bit patterns that are not non-canonical NaNs. In particular
$\enc_\varphi$ is injective.
\end{lemma}
\begin{proof}
Binary interchange formats encode every non-NaN datum, each signed zero and each infinity by
exactly one bit pattern~\cite{ieee754-2019}; decimal formats, which have several encodings per
value, are not used. Hence $\dec(\enc(v))=v$ for $v\neq\NaN$, and $\dec(c_\varphi)=\NaN$. For a
pattern $p$, $\enc(\dec(p))=p$ unless $p$ encodes a NaN other than $c_\varphi$, in which case it
is $c_\varphi$; applying $\enc\circ\dec$ again changes nothing. Injectivity follows from the
left inverse.
\end{proof}

\subsection{Nodes}

\begin{definition}[Op symbols]\label{def:opsym}
An \emph{op symbol} $f$ is an opcode with its mode byte, its operand types
$(\tau_1,\dots,\tau_{n})$ and its result types $(\rho_1,\dots,\rho_{r})$. A \emph{primitive}
op symbol has a domain $D_f\subseteq\sem{\tau_1}\times\dots\times\sem{\tau_n}$ and a meaning
$\sem{f}: D_f\to\sem{\rho_1}\times\dots\times\sem{\rho_r}$; it is \emph{total} when $D_f$ is the
whole product. A set $\mathit{Comm}$ of binary op symbols is designated commutative
(\op{ADD}, \op{MUL}, \op{AND}, \op{OR}, \op{XOR}, \op{EQ}, \op{NE}, \op{BAND}, \op{BOR},
\op{FADD}, \op{FMUL}, \op{FMIN}, \op{FMAX}), and \op{FMA} is commutative in its first two
operands. \emph{Call-like} op symbols (\op{CALL\_HASH}, \op{ASSERT\_GUARD}, \op{SIMD\_MAP},
\op{FOLD\_BOUNDED}, \op{GENERATE}) take callee operands; \emph{effect} op symbols
(\op{FILE\_READ}, \op{FILE\_WRITE}, \op{CLOCK\_READ}) take an $\Eff$ token first and return one
first.
\end{definition}

The domain of a partial primitive op is its \emph{obligation} in the specification: for
example $D_{\op{ADD.checked}}$ contains the pairs whose exact sum fits the type,
$D_{\op{INDEX}}$ the pairs $(xs,i)$ with $i<\mathrm{len}(xs)$, and $D_{\op{FCOPYSIGN}}$ the
pairs whose second operand is not $\NaN$ (\cref{sec:float}).

\begin{definition}[Nodes]\label{def:node}
A node is a tuple $N=\langle P,\Theta,R,\varepsilon,(o_0,\dots,o_{k-1}),\mathit{out},C\rangle$
where $P=(\tau_1,\dots,\tau_m)$ are the parameter types;
$\Theta=((\varphi_1,lo_1,hi_1),\dots,(\varphi_l,lo_l,hi_l))$ are the learnable parameters, each
with a float format and bounds; the \emph{inputs} of $N$ are its parameters followed by its
learnable parameters, numbered $0,\dots,m+l-1$; $R=(\rho_1,\dots,\rho_r)$ are the result types;
$\varepsilon$ is a subset of the twelve effect names (\textsf{ARENA}, \textsf{CLOCK},
\textsf{FILE\_READ}, \textsf{FILE\_WRITE}, \dots); each op is $o_i=f_i(a_{i,1},\dots,a_{i,n_i})$
with operands
\[
a ::= \mathsf{in}(j) \mid \mathsf{const}(c) \mid \mathsf{callee}(h) \mid \mathsf{inv}(h) \mid
\mathsf{op}(i',s),
\]
denoting an input, a typed constant, the hash of a callee node, an invariant clause, and the
$s$-th result of op $o_{i'}$; $\mathit{out}=(a_1,\dots,a_r)$ are the output operands; and $C$
is the contract (\cref{sec:contracts}).
\end{definition}

\begin{definition}[Well-formed nodes]\label{def:wf}
A node is well formed when:
\begin{enumerate}[label=(W\arabic*),ref=W\arabic*]
\item\label{wf:order} \emph{Order.} Every operand $\mathsf{op}(i',s)$ of $o_i$ has $i'<i$ and
  $s$ below the number of results of $o_{i'}$, and an invariant never names a port that
  depends on the fold that uses it. Hence the op graph is acyclic and every port is defined
  before use.
\item\label{wf:typing} \emph{Local typing.} Operand types match the op symbol, and output
  types match $R$.
\item\label{wf:eff} \emph{Linear effect token.} If $\varepsilon\not\subseteq\{\textsf{ARENA}\}$,
  then input $0$ and result $0$ have type $\Eff$ and no other input or result has; otherwise no
  input or result has type $\Eff$. Each $\Eff$ value (input~$0$, and the $\Eff$ result of each
  op that returns one) is used exactly once in total, by the operands of ops and by
  $\mathit{out}$.
\item\label{wf:kernel} \emph{Kernels.} The callees of \op{SIMD\_MAP}, \op{FOLD\_BOUNDED} and
  \op{GENERATE} have an empty effect set, fixed-size parameter and result types, and one
  result.
\item\label{wf:effects} \emph{Effects.} An effect op occurs only if its effect is in
  $\varepsilon$, and $\varepsilon_g\subseteq\varepsilon$ for every callee $g$.
\item\label{wf:handles} \emph{Handles.} Values of type $\Eff$ or of a handle type occur only as
  top-level inputs and results, as call arguments and results, and as effect-op operands; never
  inside a tuple, result or collection, in a kernel signature, or as a constant.
\item\label{wf:calls} \emph{Calls.} The call graph is acyclic.
\item\label{wf:capacity} \emph{Field capacity.} Every count stored in a fixed-width field of
  the binary format fits that field (\cref{cond:capacity}).
\end{enumerate}
\end{definition}

The first property is what the specification calls define-before-use; it makes every prefix
of a VG-Sym file checkable (invariant~I3). The next lemma is the formal content of the
statement that the effect token ``orders'' effects.

Kernels are pure, but the callees of \op{ASSERT\_GUARD} are ordinary nodes and may have
effects. Such a guard takes the effect token as an argument and returns its selected callee's
token, so, like a call to an effectful node, it consumes one $\Eff$ value and produces one; only
the selected callee's effects happen.

\begin{lemma}[The effect chain]\label{lem:chain}
Let $N$ be well formed with $\varepsilon\not\subseteq\{\textsf{ARENA}\}$. The ops that consume
an $\Eff$ value form a single sequence $e_1,\dots,e_p$ such that $e_1$ consumes input $0$,
$e_{q+1}$ consumes the $\Eff$ result of $e_q$, and output $0$ is the $\Eff$ result of $e_p$
(or input $0$ when $p=0$). Hence the effectful ops of $N$, including calls to effectful nodes,
are totally ordered.
\end{lemma}
\begin{proof}
Let $G$ be the directed graph whose vertices are the $\Eff$ values of $N$, with an edge
$u\to w$ when the op that produces $w$ consumes $u$. Every op that consumes an $\Eff$ value
consumes exactly one and produces exactly one (\ref{wf:typing}, \ref{wf:eff}), so every vertex
other than input $0$ has in-degree exactly $1$, and input $0$ has in-degree $0$. Every vertex
is used exactly once (\ref{wf:eff}): either by an op, giving out-degree $1$, or by
$\mathit{out}$, giving out-degree $0$; since only result $0$ has type $\Eff$, exactly one
vertex is used by $\mathit{out}$. By \ref{wf:order} $G$ is acyclic. A finite acyclic graph
whose vertices all have in- and out-degree at most $1$ is a disjoint union of paths, one per
vertex of in-degree $0$; here there is one. So $G$ is a single path from input $0$ to the
vertex used by $\mathit{out}$, and it contains every $\Eff$ value. Listing the producers of
its vertices in path order gives $e_1,\dots,e_p$.
\end{proof}

\subsection{Evaluation}

\begin{definition}[Environment]\label{def:env}
An environment consists of a callee semantics $\eta$ and an effect trace $T=(t_1,t_2,\dots)$,
the sequence of value results that the environment supplies to effect ops. $\eta$ maps each
callee hash to a partial function; for an effectful callee, the function also takes the unread
rest of the trace and returns the part it has not read.
\end{definition}

\begin{definition}[Evaluation]\label{def:eval}
Fix a well-formed node $N$, an input vector $v$ and an environment $(\eta,T)$. The partial
function $\val$ on operands is defined by recursion on the op index, which \ref{wf:order}
makes well founded: $\val(\mathsf{in}(j))=v_j$, $\val(\mathsf{const}(c))=c$, and
$\val(\mathsf{op}(i,s))=\pi_s(\mathit{Res}_i)$, where the results $\mathit{Res}_i$ of $o_i$
are, writing $\bar a$ for the values of its operands:
\begin{itemize}
\item for a primitive op, $\sem{f_i}(\bar a)$ if $\bar a\in D_{f_i}$;
\item for $\op{CALL\_HASH}(\mathsf{callee}(h),\bar a)$, the value $\eta(h)(\bar a)$;
\item for $\op{ASSERT\_GUARD}(c,\mathsf{callee}(h_1),\bar a,\mathsf{callee}(h_2))$, the value
  $\eta(h_1)(\bar a)$ if $c=\true$ and $\eta(h_2)(\bar a)$ otherwise; only the selected callee
  is applied;
\item for $\op{SIMD\_MAP}(\mathsf{callee}(h),xs,\bar b)$, the collection $ys$ of the same shape
  with $ys[\iota]=\eta(h)(xs[\iota],\bar b)$ for every index $\iota$;
\item for $\op{FOLD\_BOUNDED}(\mathsf{callee}(h),xs,\bar b,\mathit{init})$, the value
  $\mathit{acc}_n$ where $n=\mathrm{len}(xs)$, $\mathit{acc}_0=\mathit{init}$ and
  $\mathit{acc}_{q+1}=\eta(h)(\mathit{acc}_q,xs[q],\bar b)$;
\item for $\op{GENERATE}(\mathsf{callee}(h),n,\bar b)$, the sequence $ys$ with
  $ys[q]=\eta(h)(q,\bar b)$ for $q<n$;
\item for the ops on the effect chain of \cref{lem:chain}, taken in chain order: an effect op
  receives the pair $(\star,t)$, where $t$ is the next unread entry of $T$, and a call, or a
  guard whose selected callee is effectful, passes the unread rest of $T$ to that callee, which
  reads its own entries in the same way;
\end{itemize}
and $\mathit{Res}_i$ is undefined, a \emph{trap}, when an operand value is undefined, lies
outside the domain, or a required callee application is undefined. $N$ \emph{traps} on
$(v,\eta,T)$ when some $\mathit{Res}_i$ is undefined; otherwise its result is
$(\val(a_1),\dots,\val(a_r))$ for $\mathit{out}=(a_1,\dots,a_r)$.
\end{definition}

Evaluation is strict: every op is evaluated whether or not its value is used, and
\op{ASSERT\_GUARD} is the only op that applies one callee and not the other. The meaning of a
VG-Sym text is the meaning of its canonical node (\cref{sec:identity}), from which dead ops have
been removed; a dead partial op therefore never traps.

\begin{proposition}[Determinism]\label{prop:determinism}
The result of $N$ on $(v,\eta,T)$, or the fact that it traps, is uniquely determined by
$(v,\eta,T)$. Every evaluation order that respects data dependencies computes the same values.
\end{proposition}
\begin{proof}
\Cref{def:eval} defines $\val$ by well-founded recursion, so it is a function. The value of
each op depends only on the values of its operands, which precede it. The trace entries an
effect op receives are fixed by the order of the effect chain, a property of the graph that
\cref{lem:chain} makes well defined, and, by induction on the call graph, by the same property
of the callees on the chain; they do not depend on the order in which pure ops are executed.
\end{proof}

\subsection{Contracts}\label{sec:contracts}

Contract formulas are a small logic over the same values. Terms name inputs by position,
\kw{result} and \kw{result.j}, ports (in invariants only), bound variables, and typed
constants; they combine them with integer arithmetic over $\Z$ (no wrapping), float
arithmetic rounded to nearest-even, IEEE comparisons, the Boolean connectives,
indexing, projection, bounded quantification $\kw{forall}(i\ \kw{in}\ a..b)\,\{\phi\}$
(over the integers $a\le i<b$), the built-ins \kw{len}, \kw{is\_finite}, \kw{is\_nan},
\kw{abs}, \kw{min}, \kw{max}, \kw{as\_bytes}, \kw{same}, \kw{is\_ok}, \kw{ok}, \kw{err}, the
fold term $\kw{fold:}f(\mathit{init},xs,k,\bar b)$, the specification call
$\kw{spec:}s(\bar t)$ and kernel applications. On floats, \kw{==} is IEEE equality
($\NaN\neq\NaN$, $-0=+0$) and $\kw{same}(a,b)$ is identity.

\begin{definition}[Unspecified values]\label{def:unspecified}
An \emph{interpretation} $\iota$ assigns a value to every application outside its domain:
integer division or remainder by zero, out-of-range indexing, \kw{ok} of an error and \kw{err}
of a success, $\kw{spec:}s$ outside the precondition of $s$, $\kw{fold:}f$ with $k$ outside
$[0,\mathrm{len}(xs)]$, and an op outside its domain inside a kernel applied in a formula. We write
$\sem{\phi}^\iota_\sigma$ for the value of $\phi$ under an assignment $\sigma$ of inputs,
results and ports, and under $\iota$.
\end{definition}

Quantifying over every $\iota$ is how ``unspecified'' is made precise: no proof can depend on
an unspecified value, because it must hold for every choice of it.

\begin{definition}[Satisfying a contract]\label{def:satisfies}
Let $\mathit{Box}(v)$ state that each learnable input lies within its bounds. A node $N$
\emph{satisfies its contract} if for every input $v$ with $\mathit{Box}(v)$, every
interpretation $\iota$ with $\sem{\Pre}^\iota_v=\true$, every admissible trace $T$, and the
callee semantics of the nodes $N$ calls: $N$ does not trap on $(v,\eta,T)$, its result $y$
satisfies $\sem{\Post}^\iota_{v,y}=\true$, and the invariant of every \op{FOLD\_BOUNDED} holds
at every step of that fold.
\end{definition}

A trace is admissible when it respects the environment facts the specification states, such
as the monotonicity of successive clock reads along one chain.

The specification also requires contracts to be \emph{well defined}: each term that could
read an unspecified value adds an obligation saying that it does not, proven in the context
where the term occurs (under the left operand of \kw{\&\&} for its right operand, under the
negated left operand of \kw{||}, under the antecedent of \kw{==>}, under the range bounds
inside \kw{forall}). Write $\mathit{WD}(\phi)$ for the conjunction of these conditions.

\begin{lemma}[Well-defined formulas have one value]\label{lem:wd}
If $\sem{\mathit{WD}(\phi)}^\iota_\sigma=\true$ for every $\iota$, then
$\sem{\phi}^\iota_\sigma=\sem{\phi}^{\iota'}_\sigma$ for all $\iota,\iota'$.
\end{lemma}
\begin{proof}
By induction on $\phi$. Variables and constants do not depend on $\iota$. For an application
$p(t_1,\dots,t_n)$ that may leave its domain, $\mathit{WD}$ contains $\mathit{WD}(t_k)$ for each
$k$ and the domain condition of $p$; by induction the arguments have $\iota$-independent values,
and they lie in the domain, where $p$ is not interpreted by $\iota$. For $t_1\mathbin{\&\&}t_2$,
$\mathit{WD}$ is $\mathit{WD}(t_1)\wedge(t_1\Rightarrow\mathit{WD}(t_2))$: $t_1$ has an
$\iota$-independent value; if it is false, so is the conjunction for every $\iota$; if it is
true, $\mathit{WD}(t_2)$ holds for every $\iota$ and the induction hypothesis applies to $t_2$.
The cases of \kw{||} and \kw{==>} are symmetric. For $\kw{forall}(i\ \kw{in}\ a..b)\{\psi\}$,
$\mathit{WD}$ contains $\mathit{WD}(a)$, $\mathit{WD}(b)$ and $\mathit{WD}(\psi)$ for every
integer $i$ in the range, so the range and every instance have $\iota$-independent values.
\end{proof}

\subsection{Effects and capabilities}

Effect sets are ordered by inclusion; the empty set is \emph{pure} and union is the join. The
specification enforces \ref{wf:effects} at every call and counts kernels as callees.

\begin{proposition}[Effect containment]\label{prop:effects}
Let every node of a program be well formed. During any evaluation of a node $N$, including the
evaluations of its callees, every effect op that executes has its effect in $\varepsilon_N$.
\end{proposition}
\begin{proof}
By induction on the height of $N$ in the call graph, which is acyclic (\ref{wf:calls}). Effect
ops in the body of $N$ have their effect in $\varepsilon_N$ by \ref{wf:effects}. An effect op
executed inside a callee $g$ has its effect in $\varepsilon_g$ by induction, and
$\varepsilon_g\subseteq\varepsilon_N$ by \ref{wf:effects}. Kernels have empty effect sets
(\ref{wf:kernel}).
\end{proof}

\begin{proposition}[No forged capabilities]\label{prop:caps}
Let every node of a program be well formed, and let a root node be called with handles opened
by the runtime. Every handle that an effect op receives during the evaluation is one of those
handles.
\end{proposition}
\begin{proof}
We show by induction on the height of a node $N$ in the call graph that every handle value
occurring in an evaluation of $N$ is one of the handle values among the inputs of $N$. A handle
value in the body of $N$ is the value of an operand of handle type. Constants have no handle
type and no primitive or effect op has a handle result (\ref{wf:handles},
\cref{app:ops}), and handles cannot be stored in or projected out of tuples, results or
collections (\ref{wf:handles}). So the operand is an input of $N$, or a result of a call whose
handle results are, by induction, handle inputs of the callee, that is, handle arguments
supplied by $N$, which are themselves operands of $N$ of handle type. Induction on the op
index closes the argument. For the root, the inputs are the handles the runtime opened.
\end{proof}

The specification extends this confinement to the whole process: the runtime opens the
handles named by the root's manifest, then drops every other authority through Landlock and
seccomp on Linux, a generated sandbox profile on macOS, and preopened directories under WASI.
That step relies on the operating system and is outside the scope of our proofs.
